\chapter{Настроение}
\label{ch:mood}

\section{Настроение --- отдельный сцепленный модуль, а не поле графа}

Снимем путаницу сразу. Настроение не вычитывается из спектра графа памяти.
Это отдельное состояние --- низкоразмерный вектор в пространстве аффекта
$\mood \in \mathbb{R}^k$ (например, оси энергичность / раздражённость /
экзальтация / любопытство) --- отдельный модуль, слой, через который проходит
извлечённое содержание.

Путаница шла от слова «спектр». Есть два разных спектра, и их нельзя
смешивать: спектр лапласиана графа (структура памяти, глава~\ref{ch:topology})
и частотное разложение настроения во времени (быстрые и медленные компоненты
аффективного сигнала). Многослойность настроения относится ко второму и к
графу отношения не имеет.

\section{Инерция и многослойность}

У настроения должна быть инерция: оно зависит и от события, и от собственного
прошлого значения:
\begin{equation}
\mood_t = f(\mood_{t-1}, event_t), \qquad action_t = g(\mood_t, event_t).
\label{eq:moodinertia}
\end{equation}

Это протекающий интегратор \eqref{eq:leaky}: инерция --- это его
рекуррентность (состояние зависит от прошлого состояния).

\begin{openq}[Почему интегратор, а не выученная динамика]
Форму $f$ можно было бы выучить --- небольшой RNN/GRU поверх потока событий.
Мы берём фиксированный протекающий интегратор по трём причинам.
Интерпретируемость: у него ровно одна ручка на слой --- постоянная времени
$\tau_\ell$ (как быстро забывает), и её смысл читается сразу; у выученной
динамики «почему настроение поехало» --- это скрытое состояние, которое не
отладишь глазами. Отсутствие обучения и данных: RNN надо на чём-то
тренировать, а размеченных траекторий настроения Рин у нас нет и не
предвидится. Достаточность: главные наблюдаемые свойства --- инерция, возврат
к базовому уровню и разные скорости на разных горизонтах --- интегратор даёт
напрямую, одним стеком с разными $\tau_\ell$. Выученную динамику держим как
апгрейд на потом: если стенд покажет, что нужны нелинейные эффекты
(насыщение, гистерезис аффекта), которые линейный интегратор не тянет.
\end{openq}

\begin{intuition}[Интегратор и вектор --- одно и то же, с двух сторон]
Чтобы снять путаницу: «протекающий интегратор» и «вектор настроения» --- не
разные сущности. Интегратор --- это правило обновления \eqref{eq:leaky}, а его
состояние --- вектор в пространстве аффекта. Многослойность настроения --- это
стек таких интеграторов с разными $\tau_\ell$; «старший вектор» и «младший
вектор», о которых мы говорили, --- это ровно состояния медленного и быстрого
интеграторов. Актуальное настроение --- их суперпозиция. И, как отмечено в
разделе~\ref{sec:twolayers}, эти слои параллельны (все питаются одним потоком
событий, переноса между ними нет) --- в отличие от слоёв памяти, где перенос
есть.
\end{intuition}

Формально: стек интеграторов с разными $\tau_\ell$, актуальное настроение ---
их взвешенная суперпозиция:
\[
\mood_t^{(\ell)} = \left(1 - \frac{1}{\tau_\ell}\right)\mood_{t-1}^{(\ell)} +
\frac{1}{\tau_\ell}\,\mathrm{appraise}(event_t),
\qquad
\mood_t = \sum_\ell w_\ell \mood_t^{(\ell)}.
\]

Младший слой (малое $\tau$) даёт высокочастотную компоненту большой амплитуды
--- «дрожит сильнее». «Суперпозиция» здесь --- про частоты во времени, а не про
граф.

\section{Setpoint: медленный слой как аттрактор}

Медленный слой --- не просто демпфер, а точка притяжения. Добавляя к
\eqref{eq:leaky} возврат к базовому уровню $\mu$, получаем затухающую к
$\mu$ динамику:
\[
\mood_t^{(\mathrm{life})} =
\mu + (1 - \kappa)\bigl(\mood_{t-1}^{(\mathrm{life})} - \mu\bigr) + \ldots,
\qquad 0 < \kappa < 1.
\]

Здоровая динамика --- отрицательная обратная связь, тянущая настроение домой.
Патология --- когда сдвинулся сам $\mu$: депрессия в этой модели не длинная
экскурсия быстрого слоя, а уехавший аттрактор.

\section{Настроение зависит от графа, и граф --- от настроения}

Связь двусторонняя, и именно она делает настроение частью мышления.
\begin{itemize}
\item Граф $\to$ настроение. Извлечённые узлы несут валентность; она и есть
$event_t$ в \eqref{eq:moodinertia}.
\item Настроение $\to$ граф. Настроение смещает вероятности переходов при
обходе: узлы с валентностью, близкой к текущему настроению, получают больше
активации:
\begin{equation}
\mathrm{mood\_congruence}(B \mid \mood) =
\exp\!\bigl(-\gamma \lVert \mathrm{valence}(B) - \mood \rVert^2\bigr).
\label{eq:moodcong}
\end{equation}
\end{itemize}

Это конгруэнтное настроению припоминание: грустный чаще вспоминает грустное.
Настроение остаётся собственным состоянием, сцепленным сбоку, а не полем
графа.

\section{Опасность петли и три тормоза}

\begin{pitfall}[Руминация]
Конгруэнтное припоминание \eqref{eq:moodcong} плюс инерция
\eqref{eq:moodinertia} дают петлю с положительной обратной связью: грустная
вспоминает грустное $\to$ грустнеет $\to$ вспоминает ещё более грустное.
Устойчивость петли определяется её петлевым усилением $g$: при $g \ge 1$
спираль расходится. Медленный слой уменьшает эффективное $g$ --- вот его
стабилизирующая роль.
\end{pitfall}

Что удерживает живой ум? Три независимых тормоза; патология --- когда откажут
все три сразу:
\begin{enumerate}
\item Setpoint/гомеостаз. Медленный слой-аттрактор снижает петлевое усиление
ниже единицы.
\item Метакогнитивное разотождествление. Акт «почему я так часто злюсь?»
переводит из «быть-в-настроении» в «наблюдать-настроение» --- второй порядок,
и переход рвёт петлю первого порядка.
\item Экзогенное прерывание. Мир вламывается ортогональным событием;
внимание конечно, внешний вход вытесняет петлю.
\end{enumerate}

Инженерное требование: (i) член setpoint; (ii) монитор «валентность залипла
слишком долго» --- высокая автокорреляция $\mood_t$ --- с правом запустить
корректив (переопределение, переключение внимания, вызов сна); (iii)
всегда-прерываемая экзогенная значимость. Условие безопасности: три тормоза не
могут быть отключены одновременно. Псевдокод монитора ---
в приложении~\ref{ch:algorithms}.

\begin{intuition}[Почему многослойность --- это стабилизатор, а не украшение]
Пусть быстрый слой попал в негативную спираль с усилением
$g_{\text{fast}} = 1{,}1$ (расходится). Актуальное настроение --- суперпозиция, и
медленные слои вносят устойчивую компоненту с $g \approx 0$. Если их вес
$w_{\text{slow}}$ достаточен, эффективное усиление суперпозиции
$\bar{g} = w_{\text{fast}}g_{\text{fast}} + w_{\text{slow}}\cdot 0 < 1$ ---
спираль гасится за счёт «тяжёлого маховика» медленных слоёв. Так
многомасштабность оправдана функционально: это демпфер руминации, а не только
реализм.
\end{intuition}

\section{Метакогниция}
\label{sec:metacognition}

Один из трёх тормозов руминации выше --- «метакогнитивное
разотождествление» --- был назван мимоходом, но это несправедливо:
метакогниция --- самостоятельная фича всего проекта, а не побочный эффект. Ей
место в отдельном разделе.

\begin{definition}[Метакогниция]
Метакогниция --- способность Рин наблюдать и рассуждать о собственных
внутренних состояниях: о своём настроении, об уверенности в своих убеждениях, о
пробелах в памяти. Это операции второго порядка --- не «что я думаю о мире», а
«что я думаю о своём думании и своём состоянии».
\end{definition}

Метакогниция не само собой разумеется: базовая модель по умолчанию её не
имеет. Это проектируемая возможность, и у неё несколько конкретных ролей,
разбросанных по всей архитектуре, --- здесь мы их собираем.

\begin{itemize}
\item Тормоз руминации. Наблюдение «почему я так часто злюсь?» переводит из
режима «быть-в-настроении» в «наблюдать-настроение» и рвёт петлю первого
порядка (раздел выше). Технически --- это монитор автокорреляции $\mood_t$,
который сам замечает залипание и запускает корректив.
\item Эпистемическая честность. «Похоже на $X$, но $Y$ я не проверила» --- это
метакогнитивный отчёт об уверенности: Рин читает поле \texttt{conf} и зону
незнания ($|\texttt{conf}| < h$) и говорит о них, вместо ложного «я уверена»
(глава~\ref{ch:epistemics}).
\item Порождение любопытства. Замечая собственный пробел или высокую ошибку
предсказания, метакогниция ставит \texttt{need\_to\_discover} --- то есть сама
формулирует, чего Рин не знает и что стоит проверить
(главы~\ref{ch:epistemics}, \ref{ch:fusion}).
\item Самоотчёт. Способность сказать пользователю о своём состоянии («устала
от одного и того же бага», «мне сейчас не до шуток») --- это метакогниция,
считывающая настроение и мотивацию и выражающая их через Rinness
(глава~\ref{ch:voice}).
\end{itemize}

\smallskip\noindent\textbf{Субстрат метакогниции: самотелеметрия.}
Метакогниции не нужен отдельный граф. Её субстрат --- лог самонаблюдения:
кольцевой буфер временного ряда собственных состояний (вектор настроения) и
счётчиков команд (\texttt{need\_to\_*}, \texttt{fusion}, \texttt{contradiction},
\texttt{retrieval}) за часы/дни/недели. Словарь внутренних команд и состояний
конечен и дискретен, поэтому телеметрия компактна и структурирована, а
статистика по ней дешева. Это не граф и не свободный текст. Но выводы
метакогниции возвращаются в граф памяти как обычные узлы-убеждения-о-себе
(с \texttt{conf}, провенансом = телеметрия-свидетельство) --- и дальше живут по
общим правилам.

\smallskip\noindent\textbf{Детекция: дешёвый монитор чисел, а не сверх-LLM.}
Первый слой --- детекция --- набор дешёвых статистических детекторов на
буфере: скользящие среднее/медиана, тренд (наклон регрессии), частота (счёт в
окне), аномалия (z-оценка от базовой линии), автокорреляция (залипание).
Детерминированно, без LLM; срабатывают по порогу/аномалии. Пример:
«blue: медиана 0,7 > среднего 0,5 за 5 дней $\Rightarrow$ несколько всплесков
перетянули среднее, база грустная»; «\texttt{need\_to\_revise}: +60\% к
30-дневной базе». В «мозг» Рин уходит не сырой ряд (много, и LLM ненадёжно
считает медиану) и не готовый вывод, а компактная структурная сводка ---
выход детектора.

\smallskip\noindent\textbf{Аномалия --- это факт в графе, а не спецпайплайн.}
И вот основной ход, который важнее всего: сводку детектора пишем в граф памяти
как обычный факт-о-себе («за 5 дней \texttt{need\_to\_revise} +60\%»; «я в
последнее время грустная»). Уверенность высокая --- это измерение, а не
догадка. Дальше --- никакого выделенного «метакогниция-LLM» и никакого
спецпайплайна «аномалия $\to$ LLM $\to$ граф». Факт просто лежит в памяти и
подхватывается тем же общим циклом, что и любая другая мысль (retrieval,
фоновый fusion). Метакогниция здесь --- не модуль, а факты-о-себе, живущие в
графе по общим правилам.

Более того, сам факт уже работает --- вообще без LLM. Узел «я в последнее
время грустная» смещает самомодель, mood-congruent retrieval и самоотчёт
(«что-то я загрустила») через обычные механизмы. Значительная часть пользы
приходит просто от записи факта; объяснение «почему» --- опциональный более
глубокий слой, а не обязательная стадия.

\smallskip\noindent\textbf{Объяснение опционально: fusion-про-себя.}
Когда причина всё же нужна (гипотезы о том, почему аномалия), работает
fusion-про-себя (глава~\ref{ch:fusion}) --- но это тот же общий композитор, а
не прокладка: он запускается обычной salience-обработкой, когда факт-аномалия
всплывает. Гипотезы берут посылки из телеметрии: «почему больше
\texttt{need\_to\_revise}?» $\to$ H1 «много новых рассуждений», H2 «кладу
непроверенные факты», H22 «читала книгу с ошибками». Каждая --- провизорный
узел, низкая уверенность (в отличие от факта-измерения --- высокой),
провенанс из телеметрии.

Проверка --- по различающему предсказанию (научный метод): H1 $\Rightarrow$
должен вырасти и \texttt{fusion-count}; H2 $\Rightarrow$ у пересмотренных узлов
короткий интервал \texttt{commit->revision} и пропущенный
\texttt{need\_to\_discover}; H22 $\Rightarrow$ они делят один источник в
провенансе. \texttt{need\_to\_discover} кодирует, какой запрос выполнить;
curiosity/scrub его выполняет. Разрешённая гипотеза оседает в граф; если ничего
не разводит --- остаётся честное «неизвестно почему», а не выдуманная причина.

\smallskip\noindent\textbf{Eager или lazy --- из salience, а не из расписания.}
Раз аномалия --- это факт, «сразу разбираться или потом» перестаёт быть
отдельной развилкой в коде: это следствие salience самого факта. Срочная
аномалия (спираль настроения, руминация) несёт высокую значимость
$\Rightarrow$ всплывает и обрабатывается общим циклом немедленно (корректив
сейчас). Несрочная (\texttt{revise-rate} вырос) просто лежит и объясняется
лениво --- когда её достанет retrieval или придёт прямой вопрос о текущем
состоянии. Дешёвые детекторы при этом крутятся в фоне всегда (они дёшевы; их
интенсивность привязана к $\rho$-бюджету, глава~\ref{ch:walker}), а дорогой
общий LLM тратится лишь на то, что salience подняла наверх. Ни чистый
on-demand (руминация подавляет самозапрос именно тогда, когда он нужнее), ни
фиксированное расписание (расточительно) не нужны.

\smallskip\noindent\textbf{Защита от конфабуляции.}
Самообъяснение --- самый склонный к выдумыванию вид рассуждения: люди уверенно
конфабулируют причины своих настроений и поступков (классика --- Nisbett \&
Wilson, 1977). Поэтому телеметрия-лог здесь не удобство, а противоядие: он
делает самогипотезы проверяемыми против реальных данных, а не сочинёнными.
Метакогниция без телеметрии = конфабуляция; с телеметрией и различающими
проверками = фальсифицируемая.

\begin{pitfall}[Само-факты не должны стать самосбывающейся ловушкой]
Идея «аномалия --- факт в графе» сильна, но у неё есть тёмная сторона. Узел «я
в последнее время грустная» (\texttt{self\_model/mood\_summary}), показанный в
\texttt{context-packer} как обычный identity-факт, запускает петлю:
mood-congruent recall начнёт сам поддерживать это состояние. Поэтому такие
узлы: (1) диагностичны и time-bounded (с явным expiry, а не «навсегда»); (2)
пакуются в промпт как наблюдение о недавней телеметрии, а не как черта:
«recent telemetry suggests ...», а не «Rin is ...». Разница между «в последнее
время грустится» и «я грустная» --- это разница между диагнозом и приговором.
\end{pitfall}

\smallskip\noindent\textbf{Метакогниция о метакогниции (второй порядок).}
Тот же монитор можно навести на телеметрию самой метакогниции: частоту
порождаемых самогипотез, долю разрешённых, долю опровергнутых. Это
\textit{sanitizer} первого порядка: ловит analysis paralysis (много гипотез, ни
одна не разрешается $\to$ перестать копаться), ненадёжную интроспекцию
(самовыводы то и дело опровергаются $\to$ доверять самоотчёту меньше ---
прямая защита от конфабуляции) и залипшие петли. Механика та же, просто этажом
выше, поэтому дёшево. Третий порядок («думаю о том, как думаю о думании»)
даёт быстро убивающую отдачу и риск регресса --- потолок разумно ставить на
втором (аналогия со второй производной в матанализе).

\smallskip\noindent\textbf{Как это соотносится с литературой.}
Идея «метакогниция = System-2, наблюдающая System-1» и то, что мета-мысли
возвращаются в память и затем влияют на действия, --- из Tou, MacAdam, Tabor
(2024), \textit{«Metacognition is all you need?»} (поверх генеративных агентов
Park и др.). Что мы берём: самонаблюдение как отдельный модуль; мета-выводы
обратно в память; локальные модели; дешёвая модель для рутины, большая ---
для метакогниции. Где идём дальше (и где вопросительный знак в названии честно
себя выдаёт): у них детекция --- это LLM, оценивающий свой прогресс, а у нас ---
дешёвый статистический монитор, LLM только на объяснение; у них нет проверяемых
самогипотез с различающими запросами и защиты от конфабуляции; их
метакогниция целевая (задача), наша шире (настроение, эпистемика, гигиена
памяти) --- то, что они сами относят к future work. Значит «метакогниция ---
всё, что нужно»? Нет: нужен ещё субстрат самотелеметрии и дисциплина
проверяемых гипотез.

\begin{intuition}[Второй наблюдатель]
Полезно представлять метакогницию как второго наблюдателя, который смотрит не
на мир, а на первого --- на саму Рин: на её настроение, уверенность, ход мысли.
Именно этот второй наблюдатель отличает существо, которое переживает состояние,
от существа, которое знает, что переживает его, и может это назвать. Для
проекта, где цель --- личность, а не ассистент, это не украшение, а одна из
несущих опор.
\end{intuition}

\begin{openq}[Глубина и стоимость метакогниции]
Открыто, насколько «глубоко» гонять метакогницию: второй порядок («думаю о
думании») дёшев и явно полезен, но третий и выше («думаю о том, как думаю о
думании») быстро даёт убивающую отдачу и риск зацикливания на самонаблюдении.
Где ставить потолок и как выделять на метакогницию ресурсы (тот же
$\rho$-бюджет, глава~\ref{ch:walker}) --- пока не зафиксировано.
\end{openq}
